\documentclass[a4paper,12pt]{article}

% --- PAQUETES PRINCIPALES ---
\usepackage[utf8]{inputenc}
% es-noshorthands desactiva los atajos de babel para " . < >, que de otro modo
% chocan con la sintaxis de TikZ (por ejemplo >=stealth en las flechas).
\usepackage[spanish,es-tabla,es-noshorthands]{babel}
\usepackage[margin=2.5cm]{geometry}
\usepackage{graphicx}
\usepackage{amsmath,amssymb}
\usepackage{booktabs}
\usepackage{tabularx}
\usepackage{listings}
\usepackage{xcolor}
\usepackage{hyperref}
\usepackage{fancyhdr}
\usepackage{float}

% --- GRAFICOS ---
% Se dibujan con pgfplots en lugar de incluir imagenes, para que salgan
% vectoriales y usen la misma tipografia que el resto del documento.
\usepackage{pgfplots}
\pgfplotsset{compat=1.18}
\usepgfplotslibrary{groupplots}

% Un color por estructura, el mismo en todas las figuras, de modo que el arbol
% B+ sea siempre ambar y el Sequential File siempre verde azulado. La paleta se
% valido para vision con deficiencia cromatica y cada serie lleva ademas su
% propio marcador, asi que tambien se distingue impresa en blanco y negro.
\definecolor{cHeap}{HTML}{4A5BC4}
\definecolor{cSeq}{HTML}{009B84}
\definecolor{cBtree}{HTML}{B06A12}
\definecolor{cHash}{HTML}{AE3F62}
\definecolor{cGrid}{HTML}{D8D5D0}
\definecolor{cEje}{HTML}{8A857E}
\definecolor{cNota}{HTML}{5A554E}

\pgfplotsset{
    figura/.style={
        width=0.92\linewidth, height=6.4cm,
        axis line style={cEje, line width=0.4pt},
        tick align=outside, tick pos=left,
        grid=both,
        major grid style={cGrid, line width=0.3pt},
        minor grid style={cGrid!45, line width=0.2pt},
        label style={font=\small},
        tick label style={font=\footnotesize},
        title style={font=\small\bfseries, align=center},
        legend cell align=left,
        legend style={
            font=\footnotesize, draw=cGrid, line width=0.4pt,
            fill=white, rounded corners=1pt, inner sep=3pt
        },
        every axis plot/.append style={line width=1.1pt, mark size=2.1pt},
        clip marker paths=true
    }
}
\tikzset{nota/.style={font=\scriptsize, text=cNota, align=left}}

% --- ENLACES ---
\hypersetup{
    colorlinks=true,
    linkcolor=blue!70!black,
    citecolor=blue!70!black,
    urlcolor=blue!70!black
}

% --- CONFIGURACIÓN DE LISTADOS DE CÓDIGO ---
\lstset{
    basicstyle=\ttfamily\footnotesize,
    keywordstyle=\color{blue!80!black}\bfseries,
    commentstyle=\color{green!50!black}\itshape,
    stringstyle=\color{orange!80!black},
    numbers=left,
    numberstyle=\tiny\color{gray},
    stepnumber=1,
    numbersep=8pt,
    backgroundcolor=\color{gray!6},
    frame=single,
    rulecolor=\color{gray!30},
    breaklines=true,
    showstringspaces=false,
    tabsize=2,
    captionpos=b,
    literate={á}{{\'a}}1 {é}{{\'e}}1 {í}{{\'i}}1 {ó}{{\'o}}1 {ú}{{\'u}}1
             {ñ}{{\~n}}1 {Á}{{\'A}}1 {É}{{\'E}}1 {Í}{{\'I}}1 {Ó}{{\'O}}1
             {Ú}{{\'U}}1 {Ñ}{{\~N}}1 {¿}{{?`}}1 {¡}{{!`}}1
}

% --- ENCABEZADOS Y PIES DE PÁGINA ---
\pagestyle{fancy}
\fancyhf{}
\fancyhead[L]{\small Base de Datos II}
\fancyhead[R]{\small Gestor Multimodal, Entregable 1}
\fancyfoot[C]{\thepage}
\renewcommand{\headrulewidth}{0.4pt}
\setlength{\headheight}{15pt}

\begin{document}

% ==============================================================================
% PORTADA
% ==============================================================================
\begin{titlepage}
    \centering
    \vspace*{0.5cm}
    \includegraphics[width=7.5cm]{UTEC-Logo.jpg}\par\vspace{0.8cm}

    {\large \textbf{UNIVERSIDAD DE INGENIERÍA Y TECNOLOGÍA}\par}
    \vspace{0.4cm}
    {\large Departamento de Ciencia de la Computación\par}
    \vspace{1.5cm}

    {\large \textbf{BASE DE DATOS II}\par}
    \vspace{0.5cm}

    {\Huge \textbf{Informe del Entregable 1}\par}
    \vspace{0.4cm}
    {\Large \textbf{Proyecto: Gestor de Base de Datos Multimodal}\par}
    \vspace{0.3cm}
    {\large \textit{Motor de Almacenamiento en Disco, Indexación y Telemetría de E/S}\par}

    \vspace{1.2cm}

    \begin{minipage}[t]{0.52\textwidth}
        \textbf{\large Integrantes:}\\[0.3cm]
        Margiory Alvarado Chavez\\
        Francisco Calle Ochoa\\
        Mauricio Gonzalez Kremer\\
        Nicolás Llerena Silva\\
        Diana Ñañez Andres
    \end{minipage}
    \hfill
    \begin{minipage}[t]{0.42\textwidth}
        \textbf{\large Docente:}\\[0.3cm]
        Mg. Percy Lovon Ramos\\[0.8cm]
        \textbf{\large Semestre Académico:}\\[0.2cm]
        2026-2
    \end{minipage}

    \vfill
    {\large \textbf{Lima, Perú}\par}
    {\large Setiembre, 2026\par}
\end{titlepage}

\pagenumbering{roman}
\tableofcontents
\newpage

\pagenumbering{arabic}
\setcounter{page}{1}

% ==============================================================================
% SECCIONES DE LOS DEMAS BLOQUES (pendientes de sus autores)
% ==============================================================================

\section{Introducción y Arquitectura del Sistema}
% Breve descripción de la arquitectura del gestor, componentes desacoplados
% (almacenamiento en disco, indexación, motor de consultas y visualización) y
% objetivos del entregable.
El presente proyecto consiste en el diseño e implementación desde cero de un mini-gestor de bases de datos construido desde cero sobre memoria secundaria. Los motores de bases de datos modernos desacoplan sus responsabilidades para garantizar mantenibilidad, escalabilidad y un control predecible de los recursos físicos de la máquina.\\

Para lograrlo, el sistema articula cuatro componentes modulares e interconectados: en la base, la capa de almacenamiento físico y organización de archivos gestiona el archivo binario en disco mediante bloques homogéneos de tamaño fijo y estructuras de página; sobre este soporte opera la capa de indexación relacional encargada de implementar estructuras avanzadas multinivel en disco (Árbol $B^+$ y Hashing Dinámico) para optimizar búsquedas puntuales y por rango. Dichos accesos son coordinados por el motor de consultas y planificador (parser SQL) el cual procesa sentencias declarativas (\texttt{CREATE}, \texttt{INSERT}, \texttt{SELECT}) y selecciona la ruta óptima de ejecución (\textit{IndexScan}, \textit{IndexRangeScan} o \textit{SeqScan}). Finalmente, el \textbf{backend API REST y cliente web SQL} expone los servicios de consulta e interactúa con el usuario mediante una interfaz gráfica que visualiza resultados tabulares y reporta la telemetría de entrada/salida en tiempo real.\\

El objetivo principal de este primer entregable es construir y validar el núcleo de almacenamiento de bajo nivel, asegurando transferencias estrictas de bloques y proveyendo la telemetría necesaria para contrastar experimentalmente el rendimiento de las estructuras de datos.

\section{Capa de Almacenamiento Físico}

\subsection{Layout de Página y Gestión Binaria de Bloques}

Para satisfacer las restricciones de hardware y optimizar el rendimiento del subsistema de E/S, la capa de almacenamiento opera sobre archivos binarios estructurados en unidades de transferencia homogéneas de tamaño fijo $B = 4096$ bytes. Se prescinde de utilidades de serialización de alto nivel en memoria (como el módulo \texttt{pickle} de Python) y de operaciones de lectura completa del archivo mediante llamadas globales como \texttt{.read()}. Toda interacción con disco se efectúa desplazando el puntero binario a posiciones calculadas mediante $\text{offset} = \text{page\_id} \times B$ mediante la instrucción \texttt{seek(offset)} y empaquetando tipos mediante el módulo \texttt{struct} en formato binario con codificación \textit{little-endian}.

Cada página física en disco incorpora una arquitectura de Página Ranurada(\textit{Slotted-Page Architecture}) la cual permite almacenar registros manteniendo un direccionamiento indirecto robusto. El layout del bloque se divide internamente en tres segmentos contiguos:

\begin{enumerate}
    \item \textbf{Cabecera de página (\textit{Page Header}):} Ocupa los primeros 20 bytes de la página estructurados mediante el formato \texttt{<iiiii} de \texttt{struct} (5 enteros con signo de 4 bytes):
    \begin{itemize}
        \item \texttt{page\_id}: Identificador único y correlativo del bloque físico en el archivo.
        \item \texttt{record\_count}: Número actual de tuplas activas alojadas en la página.
        \item \texttt{free\_space\_offset}: Puntero que señala el inicio del espacio contiguo disponible. Inicialmente se posiciona en el byte 4096 y decrece de manera progresiva hacia el inicio conforme se van insertando registros.
        \item \texttt{next\_page\_id} / \texttt{prev\_page\_id}: Punteros de encadenamiento bidireccional entre bloques adyacentes para permitir escaneos secuenciales rápidos.
    \end{itemize}
    \item \textbf{Directorio de ranuras (\textit{Slot Directory}):} Crece de forma ascendente después del header (a partir del byte 20). Cada slot ocupa 4 bytes (\texttt{<HH}) y almacena el desplazamiento y longitud del registro correspondiente dentro del bloque.
    \item \textbf{Área de carga útil (\textit{Tuplas}):} Los datos serializados de los registros se escriben desde el final del bloque (byte 4096) hacia el centro. La inserción es aceptada solo si la adición de un nuevo slot ($4$ bytes) más la longitud de la tupla no colisiona con el puntero \texttt{free\_space\_offset}.
\end{enumerate}

Bajo este esquema, cada tupla es referenciada mediante un id lógico denominado \textbf{Identificador de registro} ($\text{RID}$) definido por el par ordenado:
\begin{equation}
    \text{RID} = (\text{page\_id}, \text{slot\_number})
\end{equation}
El uso del $\text{RID}$ permite a los índices superiores desacoplar la clave de búsqueda de la posición absoluta de bytes en disco, facilitando la reubicación interna o compactación de datos sin invalidar los punteros de los índices.

\subsection{Instrumentación de E/S (DiskCounter)}

Para medir el costo de cada operación y evaluar el comportamiento asintótico exigido en los experimentos de rendimiento, el gestor incorpora un monitor estricto denominado \texttt{DiskCounter}. Este componente instrumenta las llamadas del gestor físico y registra de manera determinista las siguientes variables de estado:

\begin{itemize}
    \item \texttt{disk\_reads}: Cantidad acumulada de bloques físicos transferidos desde el almacenamiento secundario hacia la memoria principal mediante la llamada a \texttt{read\_page(page\_id)}.
    \item \texttt{disk\_writes}: Cantidad acumulada de bloques físicos persistidos hacia el disco mediante \texttt{write\_page(page\_id)} o generados por la asignación de nuevas páginas con \texttt{allocate\_page()}.
\end{itemize}

El DiskCounter asegura que solo se contabilicen transferencias a nivel de bloque ($4096$ bytes) y no lecturas intermedias o parciales. Asimismo, expone métodos para obtener las métricas por transacción y reiniciar los contadores entre corridas experimentales, lo cual alimenta la telemetría visualizada tanto en los endpoints de la API como en el panel web de métricas.

\section{Métodos de Organización de Archivos}
\subsection{Heap File}
% Inserción O(1) con Free-List o append, Full Table Scan y eliminación.

\subsection{Sequential File}
% Estructura dual (Área Principal + Overflow), búsqueda binaria y reorganización.

\section{Capa de Indexación Relacional}
\subsection{Árbol $B^+$ en Disco}

Para la indexación avanzada, se implementó un Árbol $B^+$ multinivel que opera
estrictamente sobre memoria secundaria. La estructura fue diseñada para
parametrizar de forma dinámica la capacidad de los nodos en función del tamaño
de bloque $B$ indicado por el gestor, optimizando así las transferencias de I/O.

\textbf{Estructura Física de los Nodos:} \\
Todo nodo interno u hoja inicia con una cabecera de 16 bytes empaquetada
mediante la librería \texttt{struct} de Python (\texttt{<i i i i>}). Esta
almacena:
\begin{itemize}
    \item \texttt{is\_leaf}: Bandera booleana (entero).
    \item \texttt{num\_keys}: Cantidad actual de claves activas en el bloque.
    \item \texttt{next\_leaf} y \texttt{prev\_leaf}: Punteros de encadenamiento
    doble, fundamentales para habilitar búsquedas por rango puramente
    secuenciales a nivel de disco.
\end{itemize}

\textbf{Indexación No Agrupada (RIDs):} \\
Para evitar la duplicación innecesaria de datos y mantener un alto factor de
ramificación, los nodos hoja no almacenan los registros completos. En su lugar,
guardan tuplas de 12 bytes conformadas por \texttt{<Key, Page\_ID, Slot\_ID>}.
Esto permite que el índice devuelva el Record Identifier (RID) exacto al motor
de consultas, el cual recupera la tupla final desde el \textit{Heap} o
\textit{Sequential File}.

\textbf{Operaciones y Costos:}
\begin{itemize}
    \item \textbf{Inserción:} Se realiza mediante búsqueda descendente. Al
    superar el \texttt{MAX\_KEYS} de un bloque, se aplica el algoritmo de
    división recursiva (\textit{split}) al 50\%, promoviendo la clave separadora
    hacia el nodo padre y asignando un nuevo bloque contiguo al final del
    archivo físico.
    \item \textbf{Búsqueda Puntual:} Limitada estrictamente a la altura del
    árbol $h = O(\log_M N)$ lecturas a disco, mitigado por el uso del
    \texttt{DiskCounter} inyectado desde la capa de almacenamiento.
    \item \textbf{Búsqueda por Rango:} Minimiza severamente el I/O al realizar
    un único descenso hacia la hoja del límite inferior, y a partir de ahí,
    itera horizontalmente utilizando el puntero \texttt{next\_leaf}, evitando
    así múltiples descensos desde la raíz.
\end{itemize}

\subsection{Hashing Dinámico (Extendible Hashing)}
% Función hash en disco, directorio/buckets y división ante colisiones.

% ==============================================================================
% SECCION 5: MOTOR DE CONSULTAS, PARSER SQL Y TELEMETRIA
% ==============================================================================

\section{Motor de Consultas, Parser SQL y Telemetría}

El motor de consultas es la capa declarativa del gestor. Recibe una sentencia
SQL, decide por qué camino conviene resolverla y la ejecuta contra las
estructuras físicas, pero no abre ningún archivo por su cuenta: toda la entrada
y salida ocurre detrás de dos interfaces que las capas de almacenamiento e
indexación implementan. Esa separación fue lo que permitió que las cuatro
partes del proyecto avanzaran en paralelo y se juntaran al final sin reescribir
nada.

Una sentencia atraviesa cinco fases antes de devolver resultados. El
\textit{lexer} convierte el texto en tokens, el \textit{parser} arma el árbol
sintáctico, el \textit{binder} resuelve los nombres contra el catálogo y
convierte cada literal al tipo real de su columna, el planificador elige la
ruta de acceso, y el ejecutor recorre el plan pidiendo páginas. El binder es la
frontera entre lo sintáctico y lo semántico: antes de esa fase la condición
\texttt{id = '101'} compara un entero contra una cadena de texto, y después ya
compara dos enteros.

\subsection{Gramática Declarativa y Planificador de Ejecución}

\subsubsection{Analizador léxico y sintáctico}

El analizador léxico está escrito a mano y recorre el texto carácter por
carácter, llevando la cuenta de la línea y la columna de cada token. Ese detalle
es lo que permite que un error de escritura vuelva con la posición exacta en
lugar de un mensaje genérico, y que el editor del cliente web mueva el cursor
al punto donde el parser se detuvo:

\begin{lstlisting}[caption={Error de sintaxis con posición}]
SELECT * FORM customers WHERE "Index" = 1;

  se esperaba FROM y se encontro 'FORM' (linea 1, columna 10)
\end{lstlisting}

El analizador reconoce comentarios de línea y de bloque, cadenas con comillas
simples duplicadas para escapar, y números enteros o con exponente. También
acepta identificadores entre comillas dobles, algo que terminó siendo necesario
por el conjunto de datos que elegimos: la primera columna del archivo se llama
\texttt{Index}, que es palabra reservada, y varias otras llevan espacios en el
nombre.

El parser es de descenso recursivo y la precedencia de operadores está
codificada en el orden de las llamadas, no en una tabla aparte. La disyunción
llama a la conjunción, esta a la negación y esta última al predicado, de modo
que \texttt{OR} liga más flojo que \texttt{AND}, y \texttt{AND} más flojo que
\texttt{NOT}. La gramática soportada cubre lo que pide el enunciado más dos
extensiones que necesitamos para las pruebas:

\begin{lstlisting}[language=SQL, caption={Sentencias reconocidas por el motor}]
CREATE TABLE empleados (id INT PRIMARY KEY, nombre CHAR(30), salario FLOAT)
  USING [HEAP | SEQUENTIAL] WITH (PAGE_SIZE = 8192);
DROP TABLE empleados;

CREATE INDEX idx_id ON empleados(id) USING [BTREE | HASH];
DROP INDEX idx_id;

INSERT INTO empleados VALUES (101, 'Ada Lovelace', 5200.0);
COPY empleados FROM 'empleados.csv' WITH (HEADER, DELIMITER ',');

SELECT * FROM empleados WHERE id = 101;
SELECT nombre FROM empleados WHERE id >= 100 AND id <= 500 ORDER BY id LIMIT 10;
SELECT * FROM empleados WHERE salario BETWEEN 5000 AND 9000 AND nombre IS NOT NULL;

DELETE FROM empleados WHERE id = 101;
EXPLAIN SELECT * FROM empleados WHERE id = 101;
\end{lstlisting}

La cláusula \texttt{WITH (PAGE\_SIZE = n)} hace que el tamaño de bloque sea una
propiedad de cada tabla y no una constante del código, que es lo que el
experimento 4 necesita variar. La sentencia \texttt{COPY} carga un archivo
delimitado sin pasar por el parser, y la razón es de medición: analizar un
\texttt{INSERT} con tres mil tuplas cuesta unos 55 milisegundos solo de
analizador léxico, así que a media millón de filas lo que estaríamos midiendo
en el experimento de inserción sería la velocidad del parser y no la del disco.

\subsubsection{Sistema de tipos y formato binario}

El motor no tiene un esquema fijo. El formato binario de cada registro se
deriva de los tipos declarados en el \texttt{CREATE TABLE}, de manera que
cambiar de conjunto de datos no obliga a tocar código. Cada tipo sabe
describirse como código de \texttt{struct}, convertir un literal a su valor de
Python y cuánto espacio ocupa:

\begin{table}[H]
\centering
\caption{Correspondencia entre tipos declarados y formato físico}
\begin{tabular}{@{}lll@{}}
\toprule
\textbf{Declaración} & \textbf{Código struct} & \textbf{Bytes} \\
\midrule
\texttt{INT} & \texttt{i} & 4 \\
\texttt{BIGINT} & \texttt{q} & 8 \\
\texttt{FLOAT} & \texttt{f} & 4 \\
\texttt{DOUBLE} & \texttt{d} & 8 \\
\texttt{BOOL} & \texttt{?} & 1 \\
\texttt{CHAR(n)}, \texttt{VARCHAR(n)} & \texttt{ns} & n \\
\texttt{DATE} & \texttt{10s} & 10 \\
\bottomrule
\end{tabular}
\end{table}

Una tabla de tres columnas declaradas como \texttt{INT PRIMARY KEY},
\texttt{CHAR(30)} y \texttt{FLOAT} produce el formato \texttt{<i30sf}, que
ocupa 38 bytes. A ese cuerpo se
le antepone un mapa de nulos de un bit por columna, redondeado a bytes enteros,
porque en un registro de ancho fijo no hay otra forma de distinguir el entero
cero de la ausencia de valor. El cálculo de cuántos registros entran en una
página toma en cuenta ese mapa y la entrada del directorio de slots, cosa que al
principio no hacíamos y produjo estimaciones optimistas en cerca de un nueve
por ciento.

\subsubsection{Selección de la ruta de acceso}

El planificador parte la cláusula \texttt{WHERE} en conjunciones de primer
nivel y se queda con las que son \textit{sargables}, es decir, las que comparan
una columna contra una constante. Solo esas pueden acotar la búsqueda; una
disyunción no, porque basta que se cumpla uno de sus lados para que la fila
califique, y por eso cualquier \texttt{OR} termina resolviéndose con un escaneo
completo más un filtro encima.

Las conjunciones sargables sobre la misma columna se pliegan en un solo rango.
Dos condiciones como \texttt{id >= 100} e \texttt{id <= 500} se convierten en el
intervalo cerrado $[100, 500]$, y si el rango resulta vacío el planificador lo
detecta antes de tocar el disco. Sobre ese rango se enumeran las rutas
disponibles según lo que el catálogo tenga registrado:

\begin{table}[H]
\centering
\caption{Reglas de despacho del planificador}
\begin{tabularx}{\textwidth}{@{}lX@{}}
\toprule
\textbf{Situación} & \textbf{Ruta elegida} \\
\midrule
Igualdad con índice \texttt{HASH} o \texttt{BTREE} & \texttt{IndexScan} \\
Rango con índice \texttt{BTREE} & \texttt{IndexRangeScan} \\
Igualdad sobre la clave primaria de una tabla \texttt{SEQUENTIAL} & \texttt{SequentialSearch} \\
Rango sobre la clave primaria de una tabla \texttt{SEQUENTIAL} & \texttt{SequentialRangeScan} \\
Cualquier otro caso, incluida toda disyunción & \texttt{SeqScan} \\
\bottomrule
\end{tabularx}
\end{table}

Lo que la ruta ganadora no garantiza se vuelve a verificar en un nodo
\texttt{Filter} colocado encima, así que el plan devuelve las filas correctas
gane quien gane. Por ejemplo, en \texttt{WHERE id = 101 AND dept = 'Ventas'} con
un índice sobre \texttt{id}, el índice resuelve la primera condición y el filtro
comprueba la segunda sobre las pocas filas que llegaron.

\subsubsection{Modelo analítico de costos}

Todas las rutas se cotizan en bloques transferidos, que es exactamente lo que
mide el \texttt{DiskCounter}. Elegimos esa unidad para poder poner el costo
estimado al lado del medido y comparar ambos en la misma escala:

\begin{table}[H]
\centering
\caption{Costo estimado de cada ruta, en bloques}
\begin{tabular}{@{}ll@{}}
\toprule
\textbf{Ruta} & \textbf{Bloques} \\
\midrule
\texttt{SeqScan} & $P$ \\
\texttt{IndexScan} sobre B$^+$ & $h + f$ \\
\texttt{IndexScan} sobre hash & $1{,}2 + f$ \\
\texttt{IndexRangeScan} & $h + \lceil f / 64 \rceil + f$ \\
\texttt{SequentialSearch} & $\lceil \log_2 (P_o + 1) \rceil + 1 + P_d$ \\
\texttt{SequentialRangeScan} & $\lceil \log_2 (P_o + 1) \rceil + \lceil f / c \rceil + P_d$ \\
\bottomrule
\end{tabular}
\end{table}

donde $P$ es el número de páginas de la tabla, $h$ la altura del índice, $f$ las
filas estimadas, $c$ los registros por página, $P_o$ las páginas del área
ordenada y $P_d$ las del área de desbordamiento.

El término $+f$ de las rutas por índice es una lectura suelta por cada RID que
se alcanza. Es la hipótesis pesimista de índice no agrupado, y es justamente lo
que hace que un índice pierda contra el escaneo completo cuando el rango se
ensancha. Ese cruce no es un defecto del modelo sino el fenómeno que el
experimento 3 busca exponer, y aparece medido en la sección correspondiente.

La selectividad sale de los valores mínimo y máximo que el catálogo guarda por
columna numérica, suponiendo distribución uniforme. Es una suposición burda,
pero es la que hace que un rango angosto prefiera el índice y uno amplio
prefiera el escaneo. Cuando no hay extremos conocidos el modelo cae a las
constantes de libro de texto, un tercio para una desigualdad suelta.

Un detalle que nos costó encontrar tiene que ver con el área de desbordamiento.
Al principio el costo de las rutas ordenadas ignoraba esa área, y como no está
ordenada hay que recorrerla entera en cada búsqueda. El plan prometía once
bloques donde se leían ochocientos treinta y cuatro. Lo resolvimos guardando en
el catálogo cuántas filas esperan sin ordenar, número que el motor incrementa al
insertar y pone en cero al reorganizar, de manera que planificar sigue sin tocar
el disco. Con eso el estimado volvió a coincidir con lo medido, y además el
planificador dejó de preferir una búsqueda binaria sobre un área principal
vacía, que no ahorraba nada.

\subsubsection{Ejecución y telemetría}

El ejecutor sigue el modelo de iteradores. Cada operador es un generador que
pide filas a su hijo a medida que las necesita, así que una cláusula
\texttt{LIMIT} realmente detiene el escaneo de abajo en lugar de materializar
toda la tabla primero. El único operador que tiene que acumular es el de
ordenamiento.

Las comparaciones siguen la lógica de tres valores del estándar. Comparar contra
\texttt{NULL} no da falso sino desconocido, y una fila solo sobrevive al
\texttt{WHERE} cuando el predicado evalúa a verdadero. Por eso una condición
como \texttt{dept = 'Ventas'} no devuelve las filas donde \texttt{dept} es nulo,
y hace falta \texttt{IS NULL} para encontrarlas.

Cada consulta reporta el desglose completo: tiempo de análisis léxico y
sintáctico, tiempo de planificación, tiempo de ejecución, y los bloques leídos y
escritos que el contador registró durante toda la sentencia. El contador es una
sola instancia compartida entre tablas e índices, de modo que las cifras de una
consulta suman a través de todas las estructuras que tocó. Mantenerlo único fue
una decisión temprana; con un contador por estructura los números del informe no
habrían cuadrado.

Una consecuencia de medir con esta precisión es que descubrimos que planificar
estaba costando entrada y salida. El planificador llamaba al índice para
preguntarle su altura, y esa llamada descendía el árbol de verdad. Toda consulta
indexada medía entonces más transferencias que las que su propio plan predecía.
La altura ahora se calcula al crear o cargar el índice y se guarda; solo se
recalcula cuando el archivo creció, cosa que se verifica mirando el tamaño del
archivo, sin mover un bloque.

\subsection{Backend API REST y Cliente Web SQL}

\subsubsection{Interfaz REST}

El motor se expone con FastAPI a través de cuatro rutas:

\begin{table}[H]
\centering
\caption{Endpoints del motor de consultas}
\begin{tabularx}{\textwidth}{@{}llX@{}}
\toprule
\textbf{Método} & \textbf{Ruta} & \textbf{Función} \\
\midrule
\texttt{POST} & \texttt{/api/query} & Ejecuta una sentencia y devuelve las tuplas junto al plan y las métricas \\
\texttt{GET} & \texttt{/api/tables} & Lista las tablas con su esquema, formato binario e índices activos \\
\texttt{POST} & \texttt{/api/tables/load} & Carga masiva de un archivo delimitado \\
\texttt{POST} & \texttt{/api/tables/reorganize} & Reorganiza el área ordenada de una tabla secuencial \\
\bottomrule
\end{tabularx}
\end{table}

La respuesta de una consulta trae las columnas, las filas, el plan como árbol
anidado para poder dibujarlo, el plan en texto plano para mostrarlo tal cual, y
el bloque de métricas. Los errores salen como HTTP 400 con una forma estable que
incluye el tipo de error y, cuando es de sintaxis, la línea y la columna:

\begin{lstlisting}[caption={Respuesta de error}]
{"error": "se esperaba FROM y se encontro 'FORM' (linea 1, columna 10)",
 "kind": "SQLSyntaxError", "line": 1, "column": 10}
\end{lstlisting}

\subsubsection{Cliente web}

El cliente es una página estática sin framework, servida por nginx, que además
hace de intermediario hacia el motor. Cliente y API comparten origen, así que el
navegador no necesita permisos de origen cruzado y la consola funciona igual
servida desde cualquier máquina. La interfaz tiene las cuatro secciones que pide
el enunciado:

\begin{enumerate}
    \item \textbf{Explorador de tablas.} Muestra cada tabla con su motor de
    almacenamiento, filas, páginas, bytes por registro y registros por página,
    la lista de columnas con sus tipos y la clave primaria marcada, y los
    índices activos indicando en qué estructura quedó cada uno. Las tablas
    secuenciales traen además un botón para reorganizar, que informa cuántos
    registros se rescataron del desbordamiento y cómo cambió el número de
    páginas.

    \item \textbf{Editor SQL.} Acepta varias líneas y se ejecuta con
    \texttt{Ctrl+Enter}. El resaltado no usa un modo SQL genérico sino el
    vocabulario exacto del analizador léxico del motor, de manera que una
    palabra que el parser rechazaría se ve apagada mientras se escribe. Un
    desplegable trae las sentencias de la demostración ya escritas, con las
    comillas dobles que el conjunto de datos exige.

    \item \textbf{Visor de resultados.} Tabla paginada de cien filas por página,
    con desplazamiento horizontal, la columna de número de fila fijada al
    margen y los valores nulos marcados como tales para distinguirlos de una
    cadena vacía.

    \item \textbf{Plan de ejecución y métricas.} Dibuja el árbol del plan con el
    tipo de cada nodo, la condición que aplica y el costo que el planificador le
    asignó, y sobre él pone dos barras comparando bloques estimados contra
    bloques medidos, más el desglose de tiempos.
\end{enumerate}

Ese último panel es lo que separa esta consola de un cliente SQL cualquiera.
Poner el costo estimado al lado del medido es lo único que este motor puede
mostrar y de lo que trata el curso, así que ocupa un panel propio en vez de
esconderse detrás de un \texttt{EXPLAIN}. En la interfaz el color está reservado
para la medición: todo lo demás se construye con neutros, y los dos únicos tonos
cromáticos son ámbar para lecturas y verde azulado para escrituras.

Durante las pruebas en el navegador encontramos tres cosas que arreglar. El
atributo que oculta el paginador lo pisaba el estilo del contenedor, de modo que
aparecía aun con una sola página. El mensaje de error repetía la posición porque
el motor ya la incluye en el texto. Y la columna de número de fila reutilizaba
el estilo de los valores nulos, lo que la hacía verse en cursiva gris.
% ==============================================================================
% SECCION 6: EXPERIMENTACION Y EVALUACION DE RENDIMIENTO
% ==============================================================================

\section{Experimentación y Evaluación de Rendimiento}

Todas las mediciones de esta sección se obtuvieron contra bloques reales en
disco, con el arnés incluido en el repositorio:

\begin{lstlisting}[language=bash, caption={Ejecución de los cuatro experimentos}]
python benchmarks/experiments.py --backend disk
\end{lstlisting}

El arnés deja un archivo CSV por experimento y un resumen dentro del
directorio \texttt{benchmarks/\allowbreak results/}. Los tiempos corresponden a una corrida completa de
unos trece minutos sobre una máquina con disco de estado sólido, y las cifras de
entrada y salida son las que el \texttt{DiskCounter} registró, no estimaciones.

\subsection{Conjunto de Datos de Prueba}

Usamos el conjunto \textbf{Customers} de
Datablist\footnote{\url{https://www.datablist.com/learn/csv/download-sample-csv-files\#customers-dataset}},
que ofrece el mismo esquema en varios tamaños. Trabajamos con los archivos de
100\,000 y 500\,000 registros, por encima del mínimo que pide el enunciado. Los
archivos no están versionados porque pesan entre 17 y 83 megabytes; el
repositorio incluye un script que los descarga.

El esquema tiene doce columnas. Medimos el ancho máximo real de cada una sobre
el archivo de cien mil registros y declaramos los \texttt{CHAR} con algo de
margen sobre ese valor:

\begin{table}[H]
\centering
\caption{Esquema de la tabla de prueba}
\begin{tabular}{@{}llrr@{}}
\toprule
\textbf{Columna} & \textbf{Tipo declarado} & \textbf{Ancho máximo} & \textbf{Bytes} \\
\midrule
\texttt{Index} & \texttt{INT PRIMARY KEY} & 6 & 4 \\
\texttt{Customer Id} & \texttt{CHAR(15)} & 15 & 15 \\
\texttt{First Name} & \texttt{CHAR(16)} & 11 & 16 \\
\texttt{Last Name} & \texttt{CHAR(16)} & 11 & 16 \\
\texttt{Company} & \texttt{CHAR(40)} & 36 & 40 \\
\texttt{City} & \texttt{CHAR(28)} & 24 & 28 \\
\texttt{Country} & \texttt{CHAR(56)} & 51 & 56 \\
\texttt{Phone 1} & \texttt{CHAR(24)} & 22 & 24 \\
\texttt{Phone 2} & \texttt{CHAR(24)} & 22 & 24 \\
\texttt{Email} & \texttt{CHAR(48)} & 44 & 48 \\
\texttt{Subscription Date} & \texttt{DATE} & 10 & 10 \\
\texttt{Website} & \texttt{CHAR(44)} & 39 & 44 \\
\midrule
\multicolumn{3}{r}{\textbf{Registro completo}} & \textbf{325} \\
\bottomrule
\end{tabular}
\end{table}

Con el mapa de nulos el registro ocupa 327 bytes en la página, lo que da doce
registros por bloque de 4\,096 bytes y 8\,334 páginas para cien mil filas. La
columna \texttt{Index} es entera, única y va de uno a $N$, así que sirve como
clave primaria y es la que los índices en disco aceptan, porque tanto el árbol
B$^+$ como el hash extensible empaquetan claves enteras.

Dos particularidades del archivo obligaron a ajustar cosas del motor. La primera
columna se llama \texttt{Index}, que es palabra reservada de SQL, y otras llevan
espacios en el nombre, de modo que la tabla se declara con los nombres entre
comillas dobles. Hacerlo así tiene una ventaja práctica: como los nombres
coinciden con la cabecera del archivo, la sentencia \texttt{COPY} empareja las
columnas sola y no hace falta enumerarlas.

Para los experimentos usamos una proyección de cinco columnas del mismo archivo,
de 102 bytes por registro y treinta y ocho por página, que es la que aparece en
las tablas que siguen.

\subsection{Experimento 1: Costo de Inserción Masiva}

Cargamos lotes crecientes en cuatro configuraciones: Heap File, Sequential File
con su reorganización incluida, y Heap File manteniendo un índice, primero
árbol B$^+$ y después hash extensible.

\begin{table}[H]
\centering
\caption{Tiempo de carga en segundos}
\begin{tabular}{@{}rrrrr@{}}
\toprule
\textbf{N} & \textbf{Heap} & \textbf{Sequential} & \textbf{Heap + B$^+$} & \textbf{Heap + Hash} \\
\midrule
1\,000 & 0{,}01 & 0{,}02 & 0{,}20 & 0{,}12 \\
10\,000 & 0{,}10 & 0{,}18 & 2{,}31 & 1{,}10 \\
50\,000 & 0{,}48 & 0{,}91 & 13{,}66 & 6{,}07 \\
100\,000 & 0{,}96 & 1{,}82 & 30{,}87 & 12{,}19 \\
250\,000 & 2{,}39 & 4{,}54 & 89{,}01 & 30{,}39 \\
500\,000 & 4{,}78 & 9{,}08 & 183{,}82 & 61{,}62 \\
\bottomrule
\end{tabular}
\end{table}

\begin{table}[H]
\centering
\caption{Bloques escritos durante la carga}
\begin{tabular}{@{}rrrrr@{}}
\toprule
\textbf{N} & \textbf{Heap} & \textbf{Sequential} & \textbf{Heap + B$^+$} & \textbf{Heap + Hash} \\
\midrule
1\,000 & 27 & 63 & 1\,043 & 2\,041 \\
10\,000 & 264 & 622 & 10\,439 & 20\,522 \\
50\,000 & 1\,316 & 3\,102 & 52\,199 & 102\,344 \\
100\,000 & 2\,632 & 6\,204 & 104\,402 & 204\,685 \\
250\,000 & 6\,579 & 15\,508 & 261\,004 & 511\,195 \\
500\,000 & 13\,158 & 31\,016 & 522\,014 & 1\,023\,922 \\
\bottomrule
\end{tabular}
\end{table}

\begin{figure}[H]
\centering
\begin{tikzpicture}
\begin{axis}[figura,
    xmode=log, ymode=log,
    xlabel={Registros insertados ($N$)},
    ylabel={Tiempo de carga (s)},
    xtick={1000,10000,100000,500000},
    xticklabels={$10^3$,$10^4$,$10^5$,$5{\cdot}10^5$},
    ytick={0.01,0.1,1,10,100},
    yticklabels={0{,}01,0{,}1,1,10,100},
    legend pos=north west,
]
\addplot[cHeap,  mark=square*]   coordinates {(1000,0.0104)(10000,0.0973)(50000,0.4832)(100000,0.9627)(250000,2.391)(500000,4.775)};
\addlegendentry{Heap}
\addplot[cSeq,   mark=triangle*] coordinates {(1000,0.0189)(10000,0.1811)(50000,0.9056)(100000,1.8226)(250000,4.5442)(500000,9.0784)};
\addlegendentry{Sequential}
\addplot[cHash,  mark=diamond*]  coordinates {(1000,0.1161)(10000,1.0975)(50000,6.0682)(100000,12.1914)(250000,30.3932)(500000,61.6204)};
\addlegendentry{Heap + Hash}
\addplot[cBtree, mark=*]         coordinates {(1000,0.2022)(10000,2.3062)(50000,13.663)(100000,30.8678)(250000,89.0097)(500000,183.8226)};
\addlegendentry{Heap + B$^+$}
\end{axis}
\end{tikzpicture}
\caption{Tiempo de carga frente al volumen insertado. Ambos ejes son
logarítmicos, de modo que una recta indica crecimiento lineal: las cuatro
configuraciones escalan linealmente y lo que las separa es la constante.}
\label{fig:e1-tiempo}
\end{figure}

\begin{figure}[H]
\centering
\begin{tikzpicture}
\begin{axis}[figura,
    xmode=log, ymode=log,
    xlabel={Registros insertados ($N$)},
    ylabel={Bloques escritos},
    xtick={1000,10000,100000,500000},
    xticklabels={$10^3$,$10^4$,$10^5$,$5{\cdot}10^5$},
    ytick={10,100,1000,10000,100000,1000000},
    yticklabels={10,100,$10^3$,$10^4$,$10^5$,$10^6$},
    legend pos=north west,
]
\addplot[cHeap,  mark=square*]   coordinates {(1000,27)(10000,264)(50000,1316)(100000,2632)(250000,6579)(500000,13158)};
\addlegendentry{Heap}
\addplot[cSeq,   mark=triangle*] coordinates {(1000,63)(10000,622)(50000,3102)(100000,6204)(250000,15508)(500000,31016)};
\addlegendentry{Sequential}
\addplot[cHash,  mark=diamond*]  coordinates {(1000,2041)(10000,20522)(50000,102344)(100000,204685)(250000,511195)(500000,1023922)};
\addlegendentry{Heap + Hash}
\addplot[cBtree, mark=*]         coordinates {(1000,1043)(10000,10439)(50000,52199)(100000,104402)(250000,261004)(500000,522014)};
\addlegendentry{Heap + B$^+$}
\end{axis}
\end{tikzpicture}
\caption{Bloques escritos durante la carga. La distancia vertical constante
entre el Heap y las configuraciones con índice es la amplificación de
escritura: casi dos órdenes de magnitud, y no se cierra al crecer $N$.}
\label{fig:e1-escrituras}
\end{figure}

El resultado más claro del experimento es que mantener un índice cuesta más que
guardar los datos. Para medio millón de registros el Heap File escribe 13\,158
bloques, uno por página llena, y tarda menos de cinco segundos. Con el árbol
B$^+$ encima la misma carga tarda treinta y ocho veces más y escribe cuarenta
veces más bloques.

La explicación está en cómo escribe cada estructura. El Heap mantiene en memoria
la página que está llenando y solo la baja a disco cuando se llena, así que
escribe una vez por página. El árbol B$^+$ y el hash, en cambio, escriben su
nodo o su bucket en cada inserción, sin reutilizar entre llamadas consecutivas
el camino que acaban de recorrer. Por eso el número de escrituras del árbol es
aproximadamente una por clave más las divisiones de nodo, y el del hash es
todavía mayor porque además actualiza su cabecera cada vez.

Esa amplificación nos afectó primero en el propio Heap File. La primera versión
del adaptador reescribía el bloque completo después de cada registro, lo que
para veinte mil filas daba 20\,123 escrituras sobre 123 páginas de datos, un
factor de ciento sesenta y tres. Mantener la página de cola en memoria bajó la
cifra a 123 escrituras exactas y subió el ritmo de carga de dieciocho mil a
doscientas sesenta mil filas por segundo. La misma idea aplicada a los índices
reduciría bastante la diferencia que muestra la tabla.

El Sequential File escribe alrededor del doble que el Heap porque además de
llenar el área de desbordamiento hay que reorganizar, lo que reescribe el área
principal ordenada con factor de carga de 0{,}75.

\subsection{Experimento 2: Búsquedas Puntuales de Igualdad}

Ejecutamos mil consultas de igualdad con claves aleatorias sobre una tabla de
cien mil registros, que ocupa 2\,631 páginas, comparando las cuatro rutas de
acceso.

\begin{table}[H]
\centering
\caption{Mil búsquedas puntuales sobre $N = 100\,000$}
\begin{tabular}{@{}lrrrr@{}}
\toprule
\textbf{Ruta} & \textbf{Lecturas} & \textbf{Desv.} & \textbf{ms} & \textbf{Desv.} \\
\midrule
Full Scan (Heap) & 2\,631{,}00 & 0{,}00 & 251{,}050 & 4{,}219 \\
Búsqueda binaria (Sequential) & 11{,}85 & 1{,}48 & 0{,}799 & 0{,}104 \\
Árbol B$^+$ & 5{,}00 & 0{,}00 & 0{,}184 & 0{,}009 \\
Hash extensible & 2{,}00 & 0{,}00 & 0{,}094 & 0{,}005 \\
\bottomrule
\end{tabular}
\end{table}

\begin{figure}[H]
\centering
\begin{tikzpicture}
\begin{axis}[figura,
    % mas angosto que el resto: las etiquetas del eje Y son largas y el ancho
    % de pgfplots no las incluye
    width=0.80\linewidth, height=5.2cm,
    xmode=log,
    xlabel={Bloques leídos por consulta},
    xmin=1.2, xmax=11000,
    xtick={1,10,100,1000,10000},
    xticklabels={1,10,100,$10^3$,$10^4$},
    ytick={1,2,3,4},
    yticklabels={Hash extensible, Árbol B$^+$, Sequential, Full Scan (Heap)},
    yticklabel style={font=\footnotesize},
    ymin=0.4, ymax=4.6,
    grid=major,
    point meta=explicit symbolic,
    nodes near coords, nodes near coords align={right},
    every node near coord/.append style={font=\scriptsize, text=cNota, xshift=2pt},
]
\addplot[xcomb, cHash,  line width=0.8pt, mark=*, mark size=2.6pt] coordinates {(2,1)     [2{,}00]};
\addplot[xcomb, cBtree, line width=0.8pt, mark=*, mark size=2.6pt] coordinates {(5,2)     [5{,}00]};
\addplot[xcomb, cSeq,   line width=0.8pt, mark=*, mark size=2.6pt] coordinates {(11.85,3) [11{,}85]};
\addplot[xcomb, cHeap,  line width=0.8pt, mark=*, mark size=2.6pt] coordinates {(2631,4)  [2\,631]};
\end{axis}
\end{tikzpicture}
\caption{Promedio de bloques leídos en mil búsquedas puntuales sobre cien mil
registros. Se usa un gráfico de puntos y no de barras porque el eje es
logarítmico: en una barra el valor lo codifica la longitud, y esa proporción se
pierde al comprimir tres órdenes de magnitud, mientras que la posición del
punto sigue siendo fiel.}
\label{fig:e2}
\end{figure}

Los números siguen lo que predice la teoría. El escaneo completo lee la tabla
entera siempre, y por eso su desviación es cero y su costo no depende de dónde
esté la clave. La búsqueda binaria sobre el área ordenada necesita alrededor de
$\log_2$ del número de páginas, y su desviación de 1{,}48 refleja que algunas
claves se resuelven un sondeo antes que otras. El árbol B$^+$ gasta cinco
bloques constantes, que son la página de metadatos más el descenso por tres
niveles. El hash es el más barato con dos lecturas fijas, porque llega al bucket
a través de un directorio que vive en memoria y el número de accesos no crece
con los datos.

La diferencia práctica es de tres órdenes de magnitud: mil trescientas veces
menos lecturas con el hash que con el escaneo completo, y dos mil seiscientas
veces menos tiempo. Sobre la tabla completa de doce columnas, que ocupa 8\,334
páginas, la misma comparación da 8\,333 lecturas y 400 milisegundos contra dos
lecturas y 0{,}28 milisegundos.

\subsection{Experimento 3: Búsquedas por Rango con Selectividad Variable}

Consultamos rangos que representan el 0{,}1\,\%, 1\,\%, 5\,\%, 10\,\% y 25\,\%
de las filas, comparando el árbol B$^+$ contra el Sequential File y el escaneo
completo sobre la misma tabla de cien mil registros.

\begin{table}[H]
\centering
\caption{Bloques leídos por selectividad}
\begin{tabular}{@{}rrrrr@{}}
\toprule
\textbf{Selectividad} & \textbf{Filas} & \textbf{Árbol B$^+$} & \textbf{Sequential} & \textbf{Full Scan} \\
\midrule
0{,}1\,\% & 100 & 105 & \textbf{17} & 2\,631 \\
1\,\% & 1\,000 & 1\,010 & \textbf{48} & 2\,631 \\
5\,\% & 5\,000 & 2\,631 & \textbf{192} & 2\,631 \\
10\,\% & 10\,000 & 2\,631 & \textbf{371} & 2\,631 \\
25\,\% & 25\,000 & 2\,631 & \textbf{906} & 2\,631 \\
\bottomrule
\end{tabular}
\end{table}

\begin{figure}[H]
\centering
\begin{tikzpicture}
\begin{axis}[figura,
    height=7.6cm,
    xmode=log, ymode=log,
    xlabel={Selectividad de la consulta (\% de filas devueltas)},
    ylabel={Bloques leídos},
    xtick={0.1,1,5,10,25},
    xticklabels={0{,}1\%,1\%,5\%,10\%,25\%},
    ytick={10,100,1000,3000},
    yticklabels={10,100,$10^3$,$3{\cdot}10^3$},
    xmin=0.07, xmax=40, ymin=10, ymax=9000,
    legend pos=south east,
]
% La marca del cruce va antes que las series para quedar por debajo de ellas.
% La nota se aloja en el hueco superior izquierdo, que ninguna serie ocupa.
\draw[cNota, dotted, line width=0.6pt] (axis cs:5,10) -- (axis cs:5,3400);
\node[nota, anchor=north west] at (axis cs:0.085,8200)
  {desde aquí el árbol B$^+$ deja de convenir\\y el plan vuelve al escaneo completo};
\draw[cNota, line width=0.4pt, ->, >=stealth]
  (axis cs:1.5,4400) .. controls (axis cs:2.8,4100) .. (axis cs:4.6,3300);

\addplot[cHeap, mark=square*, dashed] coordinates {(0.1,2631)(1,2631)(5,2631)(10,2631)(25,2631)};
\addlegendentry{Full Scan (Heap)}
\addplot[cBtree, mark=*] coordinates {(0.1,105)(1,1010)(5,2631)(10,2631)(25,2631)};
\addlegendentry{Árbol B$^+$}
\addplot[cSeq, mark=triangle*] coordinates {(0.1,17)(1,48)(5,192)(10,371)(25,906)};
\addlegendentry{Sequential File}
\end{axis}
\end{tikzpicture}
\caption{Costo de un rango según la fracción de tabla que devuelve. La línea
punteada es el escaneo completo, que no depende de la selectividad. El árbol
B$^+$ sube con pendiente cercana a uno porque paga una lectura por cada RID que
alcanza, y al tocarla deja de convenir. El Sequential File, al estar agrupado,
se mantiene por debajo en todo el recorrido.}
\label{fig:e3}
\end{figure}

Aquí aparecen los dos comportamientos que el experimento busca. El primero es el
cruce del árbol B$^+$. Hasta el uno por ciento le conviene usar el índice, pero
a partir del cinco por ciento el planificador deja de elegirlo y vuelve al
escaneo completo, que es lo que explican las celdas con 2\,631 lecturas. La
razón es que el índice es secundario y cada RID que alcanza cuesta una lectura
suelta, así que traer cinco mil filas por esa vía costaría más de cinco mil
bloques cuando la tabla entera son 2\,631.

El segundo comportamiento es que el Sequential File gana en todas las
selectividades, incluso en la más angosta. Eso pasa porque está agrupado: sus
registros viven físicamente ordenados por clave, de manera que un rango son
páginas contiguas y no hay que ir a buscar cada fila por separado. Leer el
veinticinco por ciento de la tabla le cuesta 906 bloques, menos de la mitad que
un escaneo completo, mientras que el índice secundario ya había renunciado
mucho antes.

La lección es que la elección de la organización física pesa más que la del
índice cuando las consultas son por rango, y que un índice secundario solo paga
cuando la selectividad es alta.

\subsection{Experimento 4: Sensibilidad al Tamaño de Bloque}

Repetimos la carga y una búsqueda variando el tamaño de bloque, midiendo el
factor de ramificación, la altura del árbol B$^+$ y las transferencias totales.

\begin{table}[H]
\centering
\caption{Efecto del tamaño de bloque sobre la tabla}
\begin{tabular}{@{}rrrr@{}}
\toprule
\textbf{$B$ (bytes)} & \textbf{Registros por página} & \textbf{Páginas} & \textbf{Lecturas del escaneo} \\
\midrule
1\,024 & 9 & 5\,556 & 5\,555 \\
2\,048 & 19 & 2\,632 & 2\,631 \\
4\,096 & 38 & 1\,316 & 1\,315 \\
8\,192 & 77 & 650 & 649 \\
\bottomrule
\end{tabular}
\end{table}

\begin{table}[H]
\centering
\caption{Efecto del tamaño de bloque sobre el árbol B$^+$ con cien mil claves}
\begin{tabular}{@{}rrrrr@{}}
\toprule
\textbf{$B$} & \textbf{Fan-out hoja} & \textbf{Fan-out interno} & \textbf{Altura $h$} & \textbf{Páginas del índice} \\
\midrule
1\,024 & 83 & 125 & 3 & 2\,419 \\
2\,048 & 169 & 253 & 3 & 1\,187 \\
4\,096 & 339 & 509 & 3 & 592 \\
8\,192 & 681 & 1\,021 & 2 & 295 \\
\bottomrule
\end{tabular}
\end{table}

\begin{figure}[H]
\centering
\begin{tikzpicture}
\begin{axis}[figura,
    xmode=log, log basis x=2, ymode=log,
    xlabel={Tamaño de bloque $B$ (bytes)},
    ylabel={Bloques leídos},
    xtick={1024,2048,4096,8192},
    xticklabels={1\,024,2\,048,4\,096,8\,192},
    xmin=800, xmax=11000,
    ytick={1,10,100,1000,10000},
    yticklabels={1,10,100,$10^3$,$10^4$},
    ymin=1, ymax=20000,
    legend pos=north east,
]
\addplot[cHeap, mark=square*] coordinates {(1024,5555)(2048,2631)(4096,1315)(8192,649)};
\addlegendentry{Escaneo completo}
\addplot[cBtree, mark=*] coordinates {(1024,4)(2048,4)(4096,4)(8192,3)};
\addlegendentry{Búsqueda puntual con B$^+$}
\end{axis}
\end{tikzpicture}
\caption{Efecto del tamaño de bloque sobre dos accesos distintos. Duplicar $B$
reduce a la mitad el costo de un escaneo completo, porque cada transferencia
trae el doble de registros. La búsqueda puntual apenas se mueve: solo baja de
cuatro bloques a tres cuando la altura del árbol cae de tres niveles a dos.}
\label{fig:e4}
\end{figure}

El número de registros por página se duplica cada vez que se duplica el bloque,
y en consecuencia el número de páginas y las lecturas de un escaneo completo se
reducen a la mitad. La relación no es exactamente de dos porque la cabecera de
veinte bytes y el directorio de slots pesan proporcionalmente más en un bloque
chico: con mil veinticuatro bytes la cabecera se lleva casi el dos por ciento
del bloque, contra un cuarto de por ciento con ocho mil ciento noventa y dos.

En el árbol B$^+$ el efecto es más interesante porque la altura crece de forma
logarítmica y por lo tanto cambia a saltos. Entre mil veinticuatro y cuatro mil
noventa y seis bytes la altura se mantiene en tres niveles aunque el fan-out se
cuadruplica, y solo al llegar a ocho mil ciento noventa y dos baja a dos. Lo que
sí cambia de forma continua es el tamaño del índice: pasa de 2\,419 páginas a
295, ocho veces menos espacio para las mismas cien mil claves.

Bloques más grandes reducen las transferencias, pero conviene recordar que cada
transferencia mueve más bytes. La ventaja es clara cuando las consultas son
escaneos o rangos, donde todo lo que entra en el bloque se aprovecha, y se
diluye cuando son búsquedas puntuales que solo necesitan un registro de los
setenta y siete que trajo la página.

\section{Conclusiones}

El trabajo confirmó que en un gestor sobre memoria secundaria las decisiones que
importan se miden en bloques transferidos y no en operaciones de memoria. Las
cuatro rutas de acceso que implementamos resuelven la misma consulta de igualdad
con costos que van de 2\,631 lecturas a dos, y esa diferencia no se nota en la
complejidad del código sino en cuántas veces hay que ir al disco.

El resultado que menos esperábamos es el del experimento 1. Dimos por sentado
que mantener un índice tendría un costo moderado sobre la inserción, y resultó
ser el factor dominante: cargar medio millón de registros en un Heap File toma
menos de cinco segundos, y con un árbol B$^+$ encima toma tres minutos. La causa
no es la estructura sino cuándo se decide bajar una página a disco. Cuando
corregimos eso en el Heap File, que al principio reescribía el bloque completo
después de cada registro, las escrituras bajaron de 20\,123 a 123 para las
mismas veinte mil filas. Los índices siguen sin esa optimización y ahí está la
mayor parte de la diferencia que muestran las tablas.

El experimento 3 dejó la otra lección. Un índice secundario deja de convenir
apenas la consulta toca más del uno por ciento de la tabla, porque cada RID que
alcanza cuesta una lectura suelta, mientras que una organización agrupada como
el Sequential File gana en todas las selectividades por leer páginas contiguas.
Dicho de otra forma, para consultas por rango la organización física de la tabla
pesa más que cualquier índice que se le agregue encima.

Haber construido un modelo de costos en la misma unidad que mide el contador
resultó más útil de lo previsto, porque convirtió cada discrepancia entre lo
estimado y lo medido en un indicio de que algo estaba mal. Así encontramos tres
problemas que de otro modo habrían pasado desapercibidos: que el cálculo de
registros por página ignoraba el mapa de nulos y el directorio de slots, que
planificar estaba costando entrada y salida porque el planificador descendía el
árbol para preguntar su altura, y que el costo de las rutas ordenadas no tomaba
en cuenta el área de desbordamiento, lo que hacía prometer once bloques donde se
leían ochocientos treinta y cuatro.

Queda pendiente para el segundo entregable llevar el mismo criterio a los
índices, empezando por mantener en memoria el camino entre inserciones
consecutivas, y extender el árbol B$^+$ para aceptar claves no enteras y
repetidas, que hoy lo limitan a indexar claves primarias.

% ==============================================================================
% ANEXOS
% ==============================================================================
\newpage
\appendix
\section{Anexos}

\subsection{Enlaces del Proyecto y Material Complementario}

\begin{itemize}
    \item \textbf{Repositorio de código (GitHub):} \\
    \url{https://github.com/MargioUTEC/Proyecto1-BaseDatosII-20262.git} \\
    \textit{Contiene el código fuente, los scripts de datos y las instrucciones
    de ejecución en un solo paso en el archivo \texttt{README.md}.}

    \vspace{0.3cm}

    \item \textbf{Video demostrativo:} \\
    \url{https://youtu.be/enlace-al-video} \\
    \textit{Incluye la inspección de archivos binarios en disco, la ejecución de
    consultas en el cliente web con sus métricas de entrada y salida, y la
    presentación de los resultados del benchmark automatizado.}
\end{itemize}

\subsection{Instrucciones de despliegue}

El sistema levanta cuatro servicios en un solo paso:

\begin{lstlisting}[language=bash, caption={Ejecución del sistema}]
git clone https://github.com/MargioUTEC/Proyecto1-BaseDatosII-20262.git
cd Proyecto1-BaseDatosII-20262
./scripts/fetch_dataset.sh
docker compose up --build
\end{lstlisting}

La consola SQL queda en \texttt{http://localhost:8080} y es el punto de entrada.
El motor de consultas escucha en el puerto 8001, el servicio de almacenamiento
físico en el 8000 y el de hashing dinámico en el 8002.

\subsection{Reproducción de los experimentos}

\begin{lstlisting}[language=bash, caption={Arnés de experimentación}]
python benchmarks/experiments.py --backend disk
\end{lstlisting}

Deja un archivo CSV por experimento y un resumen en Markdown dentro del
directorio \texttt{benchmarks/\allowbreak results/}. La opción
\texttt{-{}-quick} corre los mismos
cuatro experimentos con volúmenes menores para una verificación rápida.

\subsection{Scripts de verificación e inspección}

\begin{table}[H]
\centering
\caption{Utilidades incluidas en el repositorio}
\begin{tabularx}{\textwidth}{@{}lX@{}}
\toprule
\textbf{Script} & \textbf{Función} \\
\midrule
\texttt{scripts/smoke.sh} & Comprueba los cuatro servicios y las invariantes del motor \\
\texttt{scripts/demo.sh} & Recorrido guiado de diez pasos, con \texttt{-{}-pausa} para narrar \\
\texttt{scripts/inspect\_page.py} & Abre un bloque binario y muestra cabecera, slots y registros \\
\texttt{scripts/sql/} & Guiones SQL para la consola o para el cliente de consola \\
\bottomrule
\end{tabularx}
\end{table}

El inspector de páginas permite verificar directamente sobre el archivo que los
datos residen en bloques binarios. Los desplazamientos del directorio de slots
crecen hacia atrás porque los registros se escriben desde el final del bloque,
que es la arquitectura de página ranurada vista en el archivo real:

\begin{lstlisting}[caption={Salida del inspector sobre la primera página}]
CABECERA DE LA PAGINA 0   (20 bytes)
  page_id              0
  record_count         12
  free_space_offset    172
  next_page_id         -1
  prev_page_id         -1
  espacio libre        104 bytes

DIRECTORIO DE SLOTS   (4 bytes por entrada)
  slot 0    offset 3769, 327 B
  slot 1    offset 3442, 327 B
  slot 2    offset 3115, 327 B
  ...
  12 vivos de 12

REGISTROS DECODIFICADOS
  RID (0, 0)
      Index                1
      Customer Id          ffeCAb7AbcB0f07
      First Name           Jared
      Country              Eritrea
\end{lstlisting}

\end{document}
